This project investigates whether an image can be culturally reinterpreted while preserving the rest of its scene. A modular pipeline captions an image, proposes a culturally relevant edit, validates the proposal, localizes the source object, segments it, edits a bounded region, and verifies saved pixels. The implementation uses SmolVLM for captioning, Qwen2.5 through Ollama for planning, Grounding DINO and SAM for localization and segmentation, and multiple research-only editing backends.

The main engineering result is strong locality and reproducibility, not a general claim of cultural correctness. The latest six-case reference-guided pilot completed all cases, independently verified zero changed pixels outside permitted and protected regions, and reproduced saved outputs deterministically. One historical croissant-to-samosa case had a direct user pass; the five new cases remain pending user review. Earlier diffusion experiments executed successfully but produced visually unconvincing food transformations. Thus, the project demonstrates a traceable and guarded research pipeline, while visual realism, independent cultural evaluation, automatic edit authorization and performance on cluttered scenes remain open problems.

